\documentclass[12pt]{article}
\usepackage{amsmath, amssymb, geometry, enumitem}
\geometry{margin=1in}
\setlength{\parindent}{0pt}
\setlength{\parskip}{0.5em}

\title{ECON 101: Macroeconomics \\ Solutions -- Quiz 4}
\author{Instructor: Muhebullah Karimzada}
\date{March 31,  2026}

\begin{document}
\maketitle


\section*{Part A: Multiple Choice Solutions}

\begin{enumerate}[label=\textbf{\arabic*.}]

\item \textbf{Correct answer: (C)}\\
A \textbf{unit of account} means money is used to quote prices and record values. Listing a car price as \$28{,}000 uses money as a common measure of value.

\item \textbf{Correct answer: (C)}\\
The simple deposit multiplier is:
\[
\text{Multiplier}=\frac{1}{r_d}=\frac{1}{0.20}=5
\]

\item \textbf{Correct answer: (B)}\\
An open-market \textbf{sale} removes reserves from the banking system. As reserves fall, banks have fewer funds to support deposits and loans, so the money supply tends to fall.

\item \textbf{Correct answer: (C)}\\
Currency held by the public is part of the narrow money supply because it is highly liquid and can be used directly for transactions.

\item \textbf{Correct answer: (C)}\\
If banks hold excess reserves, they lend out less than the maximum possible amount. As a result, the actual increase in the money supply is \textbf{smaller} than the amount predicted by the simple multiplier.

\end{enumerate}

\hrule
\vspace{1em}
\newpage
\section*{Part B: Long Numerical Problem Solution}

\textbf{6. Money Creation and the Banking System}

\textbf{Given:}
\begin{itemize}
    \item Initial deposit = \$8{,}000
    \item Desired reserve ratio \(r_d = 0.10\)
    \item Banks lend out all excess reserves
    \item The public redeposits all loan proceeds
\end{itemize}

\subsection*{(a) Desired reserves held by Bank A}

Desired reserves are calculated as:
\[
\text{Desired reserves} = r_d \times \text{deposit}
\]

Substitute the values:
\[
\text{Desired reserves} = 0.10 \times 8{,}000 = 800
\]

\[
\boxed{\text{Desired reserves} = \$800}
\]

\subsection*{(b) Amount Bank A can initially lend out}

The bank can lend out its \textbf{excess reserves}, which equal:
\[
\text{Excess reserves} = \text{deposit} - \text{desired reserves}
\]

Substitute:
\[
\text{Excess reserves} = 8{,}000 - 800 = 7{,}200
\]

\[
\boxed{\text{Initial loan} = \$7{,}200}
\]

\subsection*{(c) Simple deposit multiplier}

The simple deposit multiplier is:
\[
m = \frac{1}{r_d}
\]

Substitute:
\[
m = \frac{1}{0.10} = 10
\]

\[
\boxed{m = 10}
\]

\subsection*{(d) Maximum total increase in deposits}

The maximum total increase in deposits is:
\[
\Delta D = m \times \text{initial deposit}
\]

Substitute:
\[
\Delta D = 10 \times 8{,}000 = 80{,}000
\]

\[
\boxed{\text{Maximum total increase in deposits} = \$80{,}000}
\]

\textbf{Interpretation:} Starting from the initial \$8{,}000 deposit, repeated lending and redepositing can eventually generate total deposits of \$80{,}000 across the banking system.

\subsection*{(e) Maximum total increase in loans}

There are two equivalent ways to calculate this.

\textbf{Method 1: Total loans = Total deposits -- Total reserves}

We know total deposits eventually become:
\[
80{,}000
\]

Since the reserve ratio is 10\%, total desired reserves must be:
\[
0.10 \times 80{,}000 = 8{,}000
\]

So total loans are:
\[
80{,}000 - 8{,}000 = 72{,}000
\]

\[
\boxed{\text{Maximum total increase in loans} = \$72{,}000}
\]

\textbf{Method 2: Use the first-round logic}

The first bank lends \$7{,}200. This amount is multiplied through the banking system:
\[
7{,}200 + 6{,}480 + 5{,}832 + \cdots
\]

This geometric series sums to:
\[
72{,}000
\]

So:
\[
\boxed{\text{Maximum total increase in loans} = \$72{,}000}
\]

\subsection*{(f) Why the actual increase in deposits may be smaller}

The answer in part (d) assumes:
\begin{itemize}
    \item banks lend out \emph{all} excess reserves,
    \item households redeposit \emph{all} borrowed funds,
    \item there are no leakages from the system.
\end{itemize}

If these assumptions do not hold, the actual increase in deposits will be smaller.

\textbf{Reason 1: Banks hold excess reserves}\\
If banks choose to keep some reserves instead of lending them out, then less money is loaned and the deposit-creation process becomes weaker.

\textbf{Reason 2: Households hold cash}\\
If households keep part of the loan proceeds as currency rather than redepositing them, then less money returns to the banking system. This reduces further rounds of lending.

\textbf{Conclusion:}
\[
\boxed{\text{The actual increase in deposits will be smaller than \$80{,}000}}
\]
because both excess reserves and cash holdings reduce the multiplier process.

\hrule
\vspace{1em}

\begin{center}
\textit{End of Solution Key}
\end{center}

\end{document}